% ======================================================================
% preface-de.tex
% Copyright (c) Markus Kohm, 2008-2026
%
% This file is part of the LaTeX2e KOMA-Script bundle.
%
% This work may be distributed and/or modified under the conditions of
% the LaTeX Project Public License, version 1.3c of the license.
% The latest version of this license is in
%   http://www.latex-project.org/lppl.txt
% and version 1.3c or later is part of all distributions of LaTeX
% version 2005/12/01 or later and of this work.
%
% This work has the LPPL maintenance status "author-maintained".
%
% The Current Maintainer and author of this work is Markus Kohm.
%
% This work consists of all files listed in MANIFEST.md.
% ======================================================================

\KOMAProvidesFile{preface-de.tex}
                 [$Date: 2026-10-19 10:36:35 +0200 (Mo, 19. Okt 2026) $
                  preface to a dedicated version]

\addchap{Vorwort zu \KOMAScript~3.49.3}

Der seit langem dokumentierte Umbau bei \KOMAScript{} ist mit dieser Version
in eine entscheidende Phase eingetreten. Erstmals werden Optionen und
Anweisungen in größerem Umfang aus Klassen und Paketen entfernt. Dies betrifft
alle Optionen, vor deren Verwendung schon lange gewarnt wird.

Wer also beispielsweise trotz aller Warnungen noch immer eine seit zwanzig
Jahren veraltete Option wie \Option{noheadsepline} verwendet, wird damit in
der aktuellen Version allenfalls noch eine Warnung erzielen. Die tatsächliche
Funktion ist den aktuellen und dokumentierten Optionen vorbehalten.
Wer trotz aller Warnungen noch immer eine veraltete Anweisung verwendet,
erhält sogar eine Fehlermeldung.

Anwender sollten darauf gefasst sein, dass es ab \KOMAScript~3.50 dann auch
bei veralteten Optionen keine Warnung mehr geben wird. Auch
Rückschaltmöglichkeiten mit Option \Option{version} soll es ab dann nicht mehr
geben.

Für Anwender, die diese Übergangsphase verpasst haben, gibt es im Wiki die
Auflistung von
\href{https://sourceforge.net/p/koma-script/wiki-de/Messages/}{Fehlermeldungen
  und Warnungen}. Darüber hinaus gibt es einen Beitrag mit Hinweisen
\href{https://sourceforge.net/p/koma-script/wiki-de/HowTo_DeprecatedOptions/}{wie
  veraltete Optionen zu ersetzen sind}.

Ich bin mir bewusst, dass einigen Anwendern dieser scheinbare
Paradigmenwechsel nicht gefallen wird. Aber absolute Kompatibilität sowohl im
Ergebnis als auch im Code ist ohnehin längst eine Illusion. Mit der rasanten
Entwicklung beim \LaTeX-Kern und bei weit verbreiteten Paketen ist sie
abhängig vom Dokument längst nicht mehr zu erreichen. Immerhin sind
Anpassungen in der Regel aber leicht umsetzbar.

Für mich als Entwickler bringt die Codevereinfachung eine deutliche
Verringerung des Aufwands. Angekündigt war der Schritt ohnehin schon lange.

Einen besonderen Dank möchte ich an dieser Stelle Marei Peischl
aussprechen. Sie ist derzeit die treibende Kraft auf dem Weg zu PDF-Tagging
mit \KOMAScript. Das hat mir den Anschub gegeben, die Umbauarbeiten weiter
voran zu treiben. Wenn es nach ihr ginge, würde bereits diese Release
\Macro{DocumentMetadata} ermöglichen.  Dass das kurzfristig und hoffentlich
kurzzeitig verschoben wurde, liegt allein in meiner Verantwortung. Dasselbe
gilt für alle Fehler, die sich in dieser Release befinden.

\bigskip\noindent
Markus Kohm, Neckarhausen im Oktober 2026
\endinput

%%% Local Variables: 
%%% mode: latex
%%% TeX-master: "scrguide-de.tex"
%%% coding: utf-8
%%% ispell-local-dictionary: "de_DE"
%%% eval: (flyspell-mode 1)
%%% End: 

% LocalWords:  Dokumentationsklasse Teststruktur Benutzeranleitung Releases
% LocalWords:  Pseudolängen Buchausgaben Sommerfeldts Paketautoren Marei
% LocalWords:  Peischl Drittpaketen Drittpakete Workarounds Tagging Release
% LocalWords:  noheadsepline Rückschaltmöglichkeiten Übergangsphase
% LocalWords:  Codevereinfachung Umbauarbeiten
